\documentclass[uplatex,dvipdfmx,a4paper,11pt]{jsarticle}
\input{../../../templates/kaisetsu.tex}
\kaisetsuprint{量子力学1　第8章「ブラケット記法・スペクトル分解・並進演算子」}{2026-11-08}
%% kaisetsu-moto	problems\量子力学1\shou08\量子力学1-8.tex	kiroku\kaisetsu\量子力学1\量子力学1-8_2026-11-08_2_komento.tex
\begin{document}
\matome
\begin{saitenichiran}
  \hyougyou{\hyperlink{bt-m-072-translation}{7}}{L3・導出}{◎◎◎△◎}{B}{ヒント1を見た / (4)方針}
  \hyougyou{\hyperlink{bt-m-134-propagator}{8}}{L4・導出}{◎◎◎▲}{B}{(4)方針}
\end{saitenichiran}
\begin{daiji}
  \item 位置が確定した状態 $\ket{x'}$ は，運動量の分布 $\abs{\braket{k}{x'}}^{2}=1/(2\pi)$ が一様で，運動量はまったく定まらない。不確定性関係で $\delta x\to0$ なら $\delta p\to\infty$ であり，「運動量が一意に定まる」のではない（問題8 (4)）。
  \item 位置表示では $\matrixel{x}{\hat{p}^{n}}{\psi}=(-\ii\hbar)^{n}\psi^{(n)}(x)$ で，並進演算子 $\hat{T}(a)$ のべき級数は $\psi(x-a)$ のテイラー展開そのものである。「確かめよ」と問われたら，前の小問と別の道筋（ここでは運動量の完全性を使わない道筋）で同じ結果を出す（問題7 (4)）。
\end{daiji}
\begin{maemo}
  \item 運動量の分布を $\abs{\braket{k}{\psi}}^{2}$ で考えない：第7章 問題3 (4)（2026-11-04）で井戸のエネルギーの固有状態について運動量で何がいえるかを答えず，第7章 問題4 (1)（2026-11-07）で $\abs{\tilde{\psi}(k)}^{2}\dd k$ が運動量の確率と読める理由を書かなかった。今回は「エネルギーが一定なら運動量が一意に定まる」と，エネルギーと運動量の確定を結びつけた。「運動量の分布」と問われたら，まず $\braket{k}{\psi}$ を書き，その絶対値の2乗が $k$ によってどう変わるかを見る。
  \item 複素共役とエルミート共役の書き方：第3章 問題2（2026-10-18・2026-10-21）で複素定数の $\abs{A}^{2}$ を $A^{2}$ と書き，第5章 問題7（2026-10-30）で $(\hat{A}\psi)^{*}$ を $\hat{A}^{\dagger}\psi^{*}$ と書いた。今回も（減点はしていないが）問題7 (3) でブラ $\bra{x}\hat{T}(a)$ を $(\hat{T}(a)^{\dagger}\ket{x})^{*}$，$\ket{x-a}^{*}$ と書き，問題7 (5) の位置表示の計算で $(\dd\psi/\dd x)^{*}$ の $*$ を2行落とした（次の行で戻している）。問題6 (5) の $(\partial\tilde{\psi}/\partial k)^{2}$ は $\tilde{\psi}$ が実数なので正しいが，一般には $\abs{\partial\tilde{\psi}/\partial k}^{2}$ である。ケットからブラへは $\dagger$（エルミート共役），数には $*$（複素共役）と使い分け，$\abs{f}^{2}=f^{*}f$ の $*$ を落とさない。
\end{maemo}

\mondaihen
\level{3}
\mondaibangou{7}
\begin{mondai}{072-translation}{導出}{連続固有値のスペクトル分解と並進演算子}
区間 $(-\infty,\infty)$ の1次元系で，記号と規約は Level 2 の運動量表示の問題と同じとする。$a$ を実定数とし，並進演算子を $\hat{T}(a)=\exp(-\ii a\hat{p}/\hbar)$（べき級数で定義）とする。
  \begin{shoumon}
    \item $\hat{p}^{n}$ と $\hat{T}(a)$ が次のように書けることを示せ（和と積分の順序は入れかえてよい）。
    \[ \hat{p}^{n}=\int_{-\infty}^{\infty}(\hbar k)^{n}\ket{k}\bra{k}\dd k,\qquad \hat{T}(a)=\int_{-\infty}^{\infty}\ee^{-\ii ka}\ket{k}\bra{k}\dd k \]
    \item $\hat{T}(a)\ket{x}=\ket{x+a}$ を示せ。
    \item $\hat{T}(a)^{\dagger}=\hat{T}(-a)$ と $\hat{T}(a)^{\dagger}\hat{T}(a)=\hat{1}$ を示せ。さらに $\psi(x)=\braket{x}{\psi}$ として，$\matrixel{x}{\hat{T}(a)}{\psi}=\psi(x-a)$ を示せ。
    \item $\psi(x)$ が $x$ のまわりでテイラー展開できるとする。$\matrixel{x}{\hat{p}^{n}}{\psi}$ を使って $\hat{T}(a)$ のべき級数を直接位置表示にし，(3) の結果を確かめよ。
    \item $\ket{\psi_{a}}=\hat{T}(a)\ket{\psi}$ とする（$\ket{\psi}$ は規格化されている）。$\matrixel{\psi_{a}}{\hat{x}}{\psi_{a}}=\matrixel{\psi}{\hat{x}}{\psi}+a$ と $\matrixel{\psi_{a}}{\hat{p}}{\psi_{a}}=\matrixel{\psi}{\hat{p}}{\psi}$ を示せ。
  \end{shoumon}

\begin{kotaeran}
  \machigai{4}{△}{方針} $\matrixel{x}{\hat{p}^{n}}{\psi}=\int(\hbar k)^{n}\braket{x}{k}\braket{k}{\psi}\dd k$ とし，級数を
  \[ \int_{-\infty}^{\infty}\frac{1}{\sqrt{2\pi}}\ee^{\ii k(x-a)}\tilde{\psi}(k)\dd k=\psi(x-a) \]
  にまとめた。計算は正しいが，運動量の完全性を挟んで $\hat{p}^{n}$ を固有値 $(\hbar k)^{n}$ に置きかえたので，(1)〜(3) と同じ道筋になり，問われている「直接位置表示にする」確かめになっていない（テイラー展開の仮定も使っていない）。位置表示の $\matrixel{x}{\hat{p}^{n}}{\psi}=(-\ii\hbar)^{n}\psi^{(n)}(x)$（要点整理）を使えば，第 $n$ 項の係数が $(-\ii a/\hbar)^{n}(-\ii\hbar)^{n}/n!=(-a)^{n}/n!$ になり，級数が $\psi(x-a)$ の $x$ のまわりのテイラー展開と項ごとに一致する。
\end{kotaeran}
\begin{rikai}
  \item 位置表示では $\hat{p}$ は $-\ii\hbar\,\dd/\dd x$ なので，$\hat{T}(a)=\exp(-\ii a\hat{p}/\hbar)$ は $\exp(-a\,\dd/\dd x)$ になる。これは Level 1 の問題1 (3) の $\exp(a\hat{D})f(x)=f(x+a)$ で $a$ を $-a$ にしたものである。運動量表示と位置表示の両方で同じ $\psi(x-a)$ が出ることが「確かめ」の中身である。
  \item (5) で「これでいい？」と書いた交換関係の方法は正しく，解答と同じである。$\hat{T}(a)$ は $\hat{p}$ のべき級数なので $\comm{\hat{p}}{\hat{T}(a)}=0$ で，(3) の $\hat{T}(a)^{\dagger}\hat{T}(a)=\hat{1}$ と合わせれば $\hat{p}$ の期待値は変わらない。
\end{rikai}
\end{mondai}

\level{4}
\mondaibangou{8}
\begin{mondai}{134-propagator}{導出}{運動量の固有ケットの完全性による時間発展の演算子の位置表示（自由粒子の伝播関数）と，その性質}
1次元の自由粒子 $\hat{H}=\hat{p}^{2}/(2m)$ について，時刻 $t>0$ の時間発展の演算子の位置表示
  \[ K(x,t;x')=\matrixel{x}{\ee^{-\ii\hat{H}t/\hbar}}{x'} \]
  を考える（伝播関数という）。状態は $\ket{\psi(t)}=\ee^{-\ii\hat{H}t/\hbar}\ket{\psi(0)}$ と時間発展する（第9章）。要点整理の $\ket{k}$ の完全性と $\braket{x}{k}$ を使い，次の積分（$a$ は $\Re a\ge0$，$a\neq0$ の複素数，$\sqrt{a}$ は実部が正の分枝）を使ってよい。
  \[ \int_{-\infty}^{\infty}\ee^{-ak^{2}+bk}\dd k=\sqrt{\frac{\pi}{a}}\,\ee^{b^{2}/(4a)} \]
  \begin{shoumon}
    \item $K$ を $k$ についての積分で表し，次の形になることを示せ。
    \[ K(x,t;x')=\sqrt{\frac{m}{2\pi\ii\hbar t}}\exp\left[\frac{\ii m(x-x')^{2}}{2\hbar t}\right] \]
    \item 時刻 $0$ の波動関数 $\psi(x',0)$ から，時刻 $t$ の波動関数が $\psi(x,t)=\int K(x,t;x')\psi(x',0)\dd x'$ で得られる理由を述べよ。また，$K$ が $x$ と $t$ の関数として自由粒子の時間を含むシュレーディンガー方程式を満たすことを確かめよ。
    \item $\psi(x',0)=\ee^{\ii k_{0}x'}$（平面波）を (2) の式で時間発展させ，結果がエネルギー $\hbar^{2}k_{0}^{2}/(2m)$ の定常状態になることを確かめよ。
    \item $\abs{K}^{2}$ が $x$ によらないことを示し，時刻 $0$ に $x'$ に局在していた粒子が，どれほど短い時間 $t$ の後にも全空間に同じ確率密度で広がることを，運動量の分布の言葉で説明せよ。
  \end{shoumon}

\begin{kotaeran}
  \machigai{4}{▲}{方針} $\abs{K}^{2}=m/(2\pi\hbar t)$ で $x$ によらないことは正しい。説明に「エネルギーが一定 $\to$ 運動量が一意に定まる。不確定性原理 $\Delta x\,\Delta p\ge\hbar/2$ より，運動量が決まれば位置の確率は全空間に広がる」と書いた。時刻 $0$ の状態 $\ket{x'}$ は位置が確定した状態で，エネルギーの固有状態でも運動量の固有状態でもない。その運動量表示は $\braket{k}{x'}=\ee^{-\ii kx'}/\sqrt{2\pi}$ で，運動量の分布 $\abs{\braket{k}{x'}}^{2}=1/(2\pi)$ はすべての $k$ で同じである。つまり運動量はまったく定まっておらず，不確定性関係の向きが逆である。また，運動量が定まった状態（平面波）は時刻 $0$ からすでに全空間に広がっているので，「局在していた粒子が短い時間の後に広がる」ことの説明にもならない。正しくは，運動量の分布が一様でどれほど大きな運動量の成分も同じ重みで含まれ，大きな運動量の成分ほど速く遠くへ進むので，どれほど短い時間の後にも全空間に広がる，と説明する。
\end{kotaeran}
\begin{rikai}
  \item 位置の固有ケット $\ket{x'}$ の運動量表示 $\braket{k}{x'}$ は絶対値が $k$ によらない。位置の不確定さが $0$ の極限では運動量の不確定さは無限大で，運動量の分布は一様になる（デルタ関数のフーリエ変換が定数になることと同じ）。
  \item 自由粒子では運動量 $\hbar k$ の成分は速さ $\hbar k/m$ で進む（各成分に位相 $\ee^{-\ii\hbar k^{2}t/(2m)}$ が付く）。$\abs{k}$ に上限が無いので，どれほど短い $t>0$ でも，どれほど遠くまで届く成分がある。
  \item 「エネルギーが定まる」と「運動量が定まる」は別のことである。自由粒子のエネルギーの固有状態でも $\pm k$ の重ね合わせがありうる（第7章 問題3）。説明の問題では，まず問題の状態が何の固有状態か（または何の固有状態でもないか）を確かめる。
\end{rikai}
\end{mondai}

\kaitouhen
\level{3}
\kaisetsubangou{7}
\begin{kaisetsu}{072-translation}
\houshin $\hat{p}$ の固有ケットによる完全性を使って，$\hat{p}$ の関数を固有値 $\hbar k$ の関数に直す（離散の場合のスペクトル分解の連続版）。その後は $\braket{k}{x}$ の形を使って $\ket{x+a}$ を作り，エルミート共役で位置表示に移る。
\keisan \shoumonkai{1} $\hat{p}\ket{k}=\hbar k\ket{k}$ を $n$ 回使うと $\hat{p}^{n}\ket{k}=(\hbar k)^{n}\ket{k}$。運動量の完全性を右から掛けて
\[ \hat{p}^{n}=\hat{p}^{n}\int_{-\infty}^{\infty}\ket{k}\bra{k}\dd k=\int_{-\infty}^{\infty}(\hbar k)^{n}\ket{k}\bra{k}\dd k \]
べき級数に入れ，和と積分を入れかえると，$(-\ii a/\hbar)^{n}(\hbar k)^{n}=(-\ii ka)^{n}$ より
\[ \hat{T}(a)=\int_{-\infty}^{\infty}\Bigl[\sum_{n=0}^{\infty}\frac{(-\ii ka)^{n}}{n!}\Bigr]\ket{k}\bra{k}\dd k=\int_{-\infty}^{\infty}\ee^{-\ii ka}\ket{k}\bra{k}\dd k \]
\shoumonkai{2} $\braket{k}{x}=\braket{x}{k}^{*}=\ee^{-\ii kx}/\sqrt{2\pi}$ なので
\[ \hat{T}(a)\ket{x}=\int_{-\infty}^{\infty}\ket{k}\,\ee^{-\ii ka}\frac{\ee^{-\ii kx}}{\sqrt{2\pi}}\dd k=\int_{-\infty}^{\infty}\ket{k}\frac{\ee^{-\ii k(x+a)}}{\sqrt{2\pi}}\dd k=\int_{-\infty}^{\infty}\ket{k}\braket{k}{x+a}\dd k=\ket{x+a} \]
最後は運動量の完全性を使った。

\shoumonkai{3} $(c\ket{k}\bra{k})^{\dagger}=c^{*}\ket{k}\bra{k}$ を被積分関数に使うと
\[ \hat{T}(a)^{\dagger}=\int_{-\infty}^{\infty}\ee^{\ii ka}\ket{k}\bra{k}\dd k=\hat{T}(-a) \]
積は，$\braket{k'}{k}=\delta(k-k')$ で $k'$ の積分を実行して
\[ \hat{T}(a)^{\dagger}\hat{T}(a)=\int_{-\infty}^{\infty}\!\!\int_{-\infty}^{\infty}\ee^{\ii k'a}\ee^{-\ii ka}\ket{k'}\braket{k'}{k}\bra{k}\dd k\,\dd k'=\int_{-\infty}^{\infty}\ket{k}\bra{k}\dd k=\hat{1} \]
ブラ $\bra{x}\hat{T}(a)$ はケット $\hat{T}(a)^{\dagger}\ket{x}=\hat{T}(-a)\ket{x}=\ket{x-a}$（(2) で $a\to-a$）のエルミート共役なので $\bra{x}\hat{T}(a)=\bra{x-a}$。よって
\[ \matrixel{x}{\hat{T}(a)}{\psi}=\braket{x-a}{\psi}=\psi(x-a) \]
\shoumonkai{4} $\matrixel{x}{\hat{p}^{n}}{\psi}$ は $\psi$ に $-\ii\hbar\,\dd/\dd x$ を $n$ 回作用させたもの $(-\ii\hbar)^{n}\psi^{(n)}(x)$ である。べき級数の各項に左から $\bra{x}$ を掛けると
\[ \matrixel{x}{\hat{T}(a)}{\psi}=\sum_{n=0}^{\infty}\frac{1}{n!}\Bigl(-\frac{\ii a}{\hbar}\Bigr)^{n}(-\ii\hbar)^{n}\psi^{(n)}(x)=\sum_{n=0}^{\infty}\frac{(-a)^{n}}{n!}\psi^{(n)}(x) \]
（$(-\ii a/\hbar)(-\ii\hbar)=\ii^{2}a=-a$。）右辺は $\psi(x-a)$ を $x$ のまわりでテイラー展開したものそのものであり，(3) と一致する。

\shoumonkai{5} $\braket{x}{\psi_{a}}=\psi(x-a)$ なので，位置の完全性を挟み $y=x-a$ と置換して
\[ \matrixel{\psi_{a}}{\hat{x}}{\psi_{a}}=\int_{-\infty}^{\infty}x\abs{\psi(x-a)}^{2}\dd x=\int_{-\infty}^{\infty}(y+a)\abs{\psi(y)}^{2}\dd y=\matrixel{\psi}{\hat{x}}{\psi}+a \]
（$\ket{\psi}$ の規格化 $\int\abs{\psi(y)}^{2}\dd y=1$ を使った。）$\hat{T}(a)$ は $\hat{p}$ のべき級数なので $\hat{p}$ と交換する。(3) の $\hat{T}^{\dagger}\hat{T}=\hat{1}$ と合わせて
\[ \matrixel{\psi_{a}}{\hat{p}}{\psi_{a}}=\matrixel{\psi}{\hat{T}(a)^{\dagger}\hat{p}\hat{T}(a)}{\psi}=\matrixel{\psi}{\hat{T}(a)^{\dagger}\hat{T}(a)\hat{p}}{\psi}=\matrixel{\psi}{\hat{p}}{\psi} \]
\kotae{\shoumonkai{1} 示した \shoumonkai{2} 示した \shoumonkai{3} $\hat{T}(a)^{\dagger}=\hat{T}(-a)$，$\hat{T}^{\dagger}\hat{T}=\hat{1}$，$\matrixel{x}{\hat{T}(a)}{\psi}=\psi(x-a)$ \shoumonkai{4} 係数が $(-a)^{n}/n!$ となり $\psi(x-a)$ のテイラー展開に一致 \shoumonkai{5} $\expval{x}$ は $a$ だけ増え，$\expval{p}$ は変わらない}
\begin{ronten}
  \item \shoumonkai{1} $\hat{p}^{n}\ket{k}=(\hbar k)^{n}\ket{k}$ を述べてから運動量の完全性を右から掛け，べき級数の和と積分の順序を入れかえて指数関数にまとめることを書く。
  \item \shoumonkai{2} $\braket{k}{x}=\braket{x}{k}^{*}$ から被積分関数が $\braket{k}{x+a}$ になることを示し，運動量の完全性で $\ket{x+a}$ に戻すことを述べる。
  \item \shoumonkai{3} $\hat{T}(a)^{\dagger}$ は被積分関数の係数を複素共役にして求めることと，$\bra{x}\hat{T}(a)$ をケット $\hat{T}(a)^{\dagger}\ket{x}=\ket{x-a}$ のエルミート共役として扱うことを書く。
  \item \shoumonkai{4} べき級数の各項に左から $\bra{x}$ を掛けて $\matrixel{x}{\hat{p}^{n}}{\psi}=(-\ii\hbar)^{n}\psi^{(n)}(x)$ を使い，得た級数が $\psi(x-a)$ の $x$ のまわりのテイラー展開と項ごとに一致すること（テイラー展開できるという仮定を使うこと）を書く。
  \item \shoumonkai{5} $\hat{x}$ の期待値では置換の後に $\ket{\psi}$ の規格化を使うこと，$\hat{p}$ の期待値では $\hat{T}(a)$ が $\hat{p}$ のべき級数なので $\hat{p}$ と交換することと $\hat{T}(a)^{\dagger}\hat{T}(a)=\hat{1}$ を根拠に挙げる。
\end{ronten}
\kenzan $a=0$ で $\hat{T}=\hat{1}$ となり，すべての式が自明に成り立つ。$-\ii a\hat{p}/\hbar$ は（長さ）$\times$（運動量）/（作用）で無次元であり，指数関数の引数として整合する。運動量表示では $\braket{k}{\psi_{a}}=\ee^{-\ii ka}\tilde{\psi}(k)$ で，$\abs{\braket{k}{\psi_{a}}}^{2}=\abs{\tilde{\psi}(k)}^{2}$ となるので，運動量の分布が変わらないこととも整合する。
\tsumazuki $\hat{T}(a)$ はケット $\ket{x}$ を $\ket{x+a}$ に移すが，波動関数は $\psi(x)\to\psi(x-a)$ となる（グラフが $+a$ だけ右に動く）。$\psi(x+a)$ と書く誤りが多い。
\end{kaisetsu}

\level{4}
\kaisetsubangou{8}
\begin{kaisetsu}{134-propagator}
\houshin 運動量の固有ケットの完全性を挟んで $\hat{H}$ を固有値に置き換え，$k$ についての（複素数の係数の）ガウス積分を公式で計算する。時間発展は $\ket{x'}$ の完全性で位置表示に直す。
\keisan
\shoumonkai{1} $\hat{H}\ket{k}=(\hbar^{2}k^{2}/(2m))\ket{k}$ なので，演算子の関数として $\ee^{-\ii\hat{H}t/\hbar}\ket{k}=\ee^{-\ii\hbar k^{2}t/(2m)}\ket{k}$。完全性 $\int\ket{k}\bra{k}\dd k=\hat{1}$ を挟むと
\[ K=\int_{-\infty}^{\infty}\braket{x}{k}\,\ee^{-\ii\hbar k^{2}t/(2m)}\braket{k}{x'}\dd k=\frac{1}{2\pi}\int_{-\infty}^{\infty}\exp\left[-\frac{\ii\hbar t}{2m}k^{2}+\ii(x-x')k\right]\dd k \]
公式で $a=\ii\hbar t/(2m)$（$\Re a=0$），$b=\ii(x-x')$ とすると
\[ \frac{b^{2}}{4a}=\frac{-(x-x')^{2}}{4\ii\hbar t/(2m)}=\frac{\ii m(x-x')^{2}}{2\hbar t},\qquad \frac{1}{2\pi}\sqrt{\frac{\pi}{a}}=\frac{1}{2\pi}\sqrt{\frac{2\pi m}{\ii\hbar t}}=\sqrt{\frac{m}{2\pi\ii\hbar t}} \]
よって示された形になる。
\shoumonkai{2} $\ket{\psi(t)}=\ee^{-\ii\hat{H}t/\hbar}\ket{\psi(0)}$ に左から $\bra{x}$ をかけ，$\ket{\psi(0)}$ の前に完全性 $\int\ket{x'}\bra{x'}\dd x'=\hat{1}$ を挟むと
\[ \psi(x,t)=\int_{-\infty}^{\infty}\matrixel{x}{\ee^{-\ii\hat{H}t/\hbar}}{x'}\braket{x'}{\psi(0)}\dd x'=\int_{-\infty}^{\infty}K(x,t;x')\psi(x',0)\dd x' \]
次に $D=x-x'$，$K=Ct^{-1/2}\ee^{\ii mD^{2}/(2\hbar t)}$（$C$ は定数）として微分すると
\begin{align*}
  \ii\hbar\pdv{K}{t}&=\ii\hbar K\left(-\frac{1}{2t}-\frac{\ii mD^{2}}{2\hbar t^{2}}\right)=K\left(-\frac{\ii\hbar}{2t}+\frac{mD^{2}}{2t^{2}}\right),\\
  -\frac{\hbar^{2}}{2m}\pdv[2]{K}{x}&=-\frac{\hbar^{2}}{2m}K\left[\left(\frac{\ii mD}{\hbar t}\right)^{2}+\frac{\ii m}{\hbar t}\right]=K\left(\frac{mD^{2}}{2t^{2}}-\frac{\ii\hbar}{2t}\right)
\end{align*}
で両者は等しく，$K$ は自由粒子のシュレーディンガー方程式を満たす。
\shoumonkai{3} (2) の式に代入し，公式を $x'$ について使う。$x'$ の2次の係数は $a=-\ii m/(2\hbar t)$ で $\Re a=0$，1次の係数は $b=-\ii mx/(\hbar t)+\ii k_{0}$，定数項は $\ii mx^{2}/(2\hbar t)$ である。
\[ \frac{b^{2}}{4a}=\frac{-(k_{0}-mx/(\hbar t))^{2}}{-2\ii m/(\hbar t)}=-\frac{\ii\hbar t}{2m}\left(k_{0}-\frac{mx}{\hbar t}\right)^{2}=-\frac{\ii\hbar k_{0}^{2}t}{2m}+\ii k_{0}x-\frac{\ii mx^{2}}{2\hbar t} \]
最後の項は定数項と打ち消し合い，係数は $\sqrt{m/(2\pi\ii\hbar t)}\sqrt{\pi/a}=\sqrt{m/(2\pi\ii\hbar t)}\sqrt{2\pi\hbar t/(-\ii m)}=1$（$\ii(-\ii)=1$）である。よって
\[ \psi(x,t)=\exp\left[\ii\left(k_{0}x-\frac{\hbar k_{0}^{2}}{2m}t\right)\right] \]
で，$\ee^{\ii k_{0}x}$ に $\ee^{-\ii Et/\hbar}$（$E=\hbar^{2}k_{0}^{2}/(2m)$）がかかった定常状態になる。これは $\ket{k_{0}}$ に $\ee^{-\ii\hat{H}t/\hbar}$ を直接作用させた結果と一致する。
\shoumonkai{4} $\abs{\ee^{\ii\theta}}=1$（$\theta$ は実数）なので $\abs{K}^{2}=m/(2\pi\hbar t)$ で，$x$ によらない。時刻 $0$ の状態 $\ket{x'}$ は位置が確定した状態で，運動量表示は $\braket{k}{x'}=\ee^{-\ii kx'}/\sqrt{2\pi}$，すなわち運動量の分布は $\abs{\braket{k}{x'}}^{2}=1/(2\pi)$ で一様である。どれほど大きな運動量の成分も同じ重みで含まれ，運動量が大きい成分ほど速く遠くへ進むので，どれほど短い時間の後にも確率密度は全空間に一様に広がる。（位置の不確定さが $0$ の極限では，運動量の不確定さは無限大になる。）
\kotae{\shoumonkai{1} $K=(1/2\pi)\int\exp\bigl[\ii k(x-x')-\ii\hbar tk^{2}/(2m)\bigr]\dd k=\sqrt{m/(2\pi\ii\hbar t)}\,\exp\bigl[\ii m(x-x')^{2}/(2\hbar t)\bigr]$ \shoumonkai{2} 完全性 $\int\ket{x'}\bra{x'}\dd x'=\hat{1}$ を挟む。$K$ はシュレーディンガー方程式を満たす（計算は本文） \shoumonkai{3} $\psi(x,t)=\exp[\ii(k_{0}x-\hbar k_{0}^{2}t/(2m))]$ \shoumonkai{4} $\abs{K}^{2}=m/(2\pi\hbar t)$。位置が確定した状態は運動量の分布が一様で，大きな運動量の成分がいくらでも速く広がるため}
\begin{ronten}
  \item \shoumonkai{1} 完全性を挟んで $\hat{H}$ を固有値に置き換えたことと，公式を $\Re a=0$ の場合に使ったこと（係数の平方根の分枝）を書く。
  \item \shoumonkai{2} $\ket{\psi(t)}=\ee^{-\ii\hat{H}t/\hbar}\ket{\psi(0)}$ と $\ket{x'}$ の完全性から積分の式を導く。
  \item \shoumonkai{4} 位置が確定した状態の運動量の分布が一様であることを根拠にする。
\end{ronten}
\kenzan 次元：$\sqrt{m/(\hbar t)}$ は $\sqrt{\text{質量}/(\text{エネルギー}\times\text{時間}^{2})}=1/\text{長さ}$ で，$K$ は $\delta(x-x')$ と同じ 1／長さの次元をもつ。指数 $m(x-x')^{2}/(\hbar t)$ は無次元。(3) の結果は $t=0$ で $\ee^{\ii k_{0}x}$ に戻る。
\tsumazuki 公式の $\sqrt{\pi/a}$ で，$a=\ii\hbar t/(2m)$ の平方根は $\ee^{\ii\pi/4}\sqrt{\hbar t/(2m)}$（実部が正の分枝）をとる。$\sqrt{\ii}$ の分枝を誤ると全体の位相がずれる。
\end{kaisetsu}

\end{document}
